\documentclass[10pt]{article}

\usepackage[utf8]{inputenc}

\usepackage[T1]{fontenc}

\usepackage{amsmath}

\usepackage{amsfonts}

\usepackage{amssymb}

\usepackage[version=4]{mhchem}

\usepackage{stmaryrd}

\usepackage{graphicx}

\usepackage[export]{adjustbox}

\graphicspath{ {./images/} }



\begin{document}

Особые точки и их классификация. Рассмотрим однородные уравнения, соответствуюшие (1) и (2):



\[

\begin{gathered}

w^{\prime}(z)=A(z) w \\

w^{(n)}+a_{1}(z) w^{(n-1)}+\ldots+a_{n}(z) w=0 .

\end{gathered}

\]



Точка $z_{0}$ называется о с обой точкой системы (3) или уравнения (4), если она является особой точкой для одного из элементов $a_{i k}(z)$ (коэффициентов $a_{i}(z)$ ). Пусть $z_{0}-$ полюс, тогда система (3) имеет фундаментальную матрицу $W(z)$ вида



\[

W(z)=\Phi(z)\left(z-z_{0}\right) P

\]



где $P$ - постоянная матрица, матрица-функция $\Phi(z)$ разлагается в ряд Лорана



\[

\Phi(z)=\sum_{-\infty}^{\infty} \varphi_{k}\left(z-z_{0}\right)^{k}

\]



сходящийся в нек-ром колыце вида



\[

0<\left|z-z_{0}\right|<R \text {. }

\]



Особая точка $z_{0}$ называется регуля рной, если ряд Лорана для $\Phi(z)$ содержит конечное число отрицательных степеней $z-z_{0}$, и и р регуля р ной в противном случае. Это косвенная классификация: она дается в терминах свойств решений, а не коэффициентов системы. Аналогично классифицируются особые точки уравнения (4). Бесконечно удаленная точка $z=\infty$ называется о с о б о й точ ко й системы (3), если точка $t=0$ - особая для системы



\[

w_{t}^{\prime}=-t^{-2} A\left(t^{-1}\right) w

\]



полученной из (3) заменой переменного $z=1 / t$; аналогично для уравнения (4).



Регуля рны е особы е точки - наиболее простой и хорошо изученный тип особых точек. Точка $z_{0}$ является регулярной особой точкой уравнения (4) тогда и только тогда, когда



\[

a_{i}(z)=\left(z-z_{0}\right)^{-i} p_{i}(z)

\]



где функции $p_{i}(z)$ аналитичны в точке $z_{0}$. Точка $z=\infty$ является регулярной особой точкой уравнения (4) тогда и только тогда, когда $a_{i}(z)=z^{-i} q_{i}(z)$, где функции $q_{i}(z)$ аналитичны в точке $z=\infty$. Определяющее уравнение в регулярной особой точке $z_{0}$ имеет вид:



\[

\begin{gathered}

\rho(\rho-1) \ldots(\rho-n+1)+p_{1}\left(z_{0}\right) \rho(\rho-1) \ldots(\rho-n+2)+\ldots \\

\ldots+p_{n}\left(z_{0}\right)=0,

\end{gathered}

\]



его корни называются характе ристическими пока з а т е л я м и в точке $z_{0}$. Если ни одна из разностей $\rho_{i}-\rho_{k}$, $i \neq k$, не есть целое число, то уравнение (4) имеет следуюшую фундаментальную систему решений:



\[

w_{i}(z)=\left(z-z_{0}\right)^{\rho_{i}} \varphi_{i}(z), \varphi_{i}\left(z_{0}\right)=1,1 \leqslant i \leqslant n,

\]



где функции $\varphi_{i}(z)$ аналитичны в точке $z_{0}$. Если среди этих разностей есть целые числа, то решения могут содержать целые степени логарифма $\ln \left(z-z_{0}\right)$.



Уравнение 2-го порядка с регулярной особой точкой $z_{0}$ имеет вид:



\[

w^{\prime \prime}+\left(z-z_{0}\right)^{-1} p_{1}(z) w^{\prime}+\left(z-z_{0}\right)^{-2} p_{2}(z) w=0,

\]



где функции $p_{1}(z), p_{2}(z)$ аналитичны в точке $z_{0}$; определяюшее уравнение таково:



\[

\rho(\rho-1)+\rho p_{1}\left(z_{0}\right)+p_{2}\left(z_{0}\right)=0 .

\]



Если $\rho_{1}-\rho_{2}-$ нецелое число, где $\rho_{i}-$ характеристич. показатели, то уравнение (5) имеет фундаментальную систему решений



\[

w_{i}(z)=\left(z-z_{0}\right)^{\rho_{i}} \varphi_{i}(z)

\]



где функции $\varphi_{i}(z)$ аналитичны в точке $z_{0}, \varphi_{i}\left(z_{0}\right)=1$. Если $\rho_{1}-\rho_{2}$ есть целое неотрицательное число, то уравнение (5) имеет фундаментальную систему решений



\[

\begin{aligned}

& w_{1}(z)=\left(z-z_{0}\right)^{\rho_{1} \varphi_{1}(z)},

\end{aligned}

\]



\begin{center}

\includegraphics[max width=\textwidth]{2023_07_08_5f876aeda83ce1d68ee8g-1}

\end{center}



где $\theta$ - постоянная, функции $\varphi_{i}(z)$ анӑлитичны в точке $z_{0}$, $\varphi_{i}\left(z_{0}\right)=1$.



Примеры: уравнение Эйри: $w^{\prime \prime}-z w=0, z=\infty-$ иррегулярная особая точка; уравнение Бе с с ля: $z^{2} w^{\prime \prime}+z w^{\prime}+\left(z^{2}-\gamma^{2}\right) w=0, z=0-$ регулярная, $z=\infty-$ иррегулярная особая точка; гиперге ометрич. уравнение: $z(1-z) w^{\prime \prime}+[\gamma-(\alpha+\beta+1) z] w^{\prime}-\alpha \beta w=0$ имеет регулярные особые точки: $0,1, \infty$.



Уравнением класса Фукса называется уравнение (4), все особые точки к-рого на римановой сфере являются регулярными. Известен обший вид таких уравнений. Все основные Д.у. 2-го порядка, возникающие в задачах математич. физики, можно получить из уравнения с пятью регулярными независимыми особыми точками; при этом разности характеристич. показателей в каждой особой точке равны $1 / 2$.



Точка $z_{0}$ является регулярной особой точкой системы (3), если



\[

A(z)=\left(z-z_{0}\right)^{-1} B(z),

\]



где матрица-функшия $B(z)$ аналитична в точке $z_{0}, B\left(z_{0}\right) \neq 0$. Если все разности $\rho_{i}-\rho_{k}, i \neq k$, где $\rho_{i}-$ собственные значения матрицы $B\left(z_{0}\right)$, не являются целыми числами, то система (3) имеет фундаментальную матрицу вида



\[

W(z)=\Phi(z)\left(z-z_{0}\right)^{P}

\]



где $P$ - диагональная матрица с элементами $\rho_{1}, \rho_{2}, \ldots, \rho_{n}$, матрица-функция $\Phi(z)$ аналитична в точке $z_{0}$ и невырождена. Если среди этих разностей есть целые числа, то фундаментальная матрица содержит целые степени $\ln \left(z-z_{0}\right)$. Неизвестны необходимые и достаточные условия того, что $z_{0}-$ регулярная особая точка системы (3). Система



\[

w^{\prime}=w \sum_{k=0}^{m} A_{k}\left(z-a_{k}\right)^{-1}

\]



где $a_{k}$ - различные комплексные числа, $A_{k}$ - постоянные ненулевые матрицы порядка $n \times n$ и $A_{1}+\ldots+A_{m} \neq 0$, является системой класса Фукса и имеет регулярные особые точки $a_{1}, a_{2}, \ldots, a_{m}, \infty$.



Иррегулярны е особые точки. Пусть в системе (3)



\[

A(z)=z^{r} \sum_{k=0}^{\infty} A_{k} z^{-k}, \quad A_{0} \neq 0

\]



где $r \geqslant 0$ - целое, ряд сходится при $|z|>R$, тогда $z=\infty$ есть иррегулярная особая точка, и система имеет фундаментальную матрицу вида



\[

W(z)=S(z) \exp Q(z)

\]



где $Q(z)$ - диагональная матрица, элементы к-рой являются многочленами от $z^{1 / n}, n>0-$ целое:



\[

q_{i i}(z)=q_{i 0} z^{l_{i} / n}+q_{i 1} z^{\left(l_{i}-1\right) / n}+\ldots+q_{i, l_{i}-1} z^{1 / n} .

\]



Өлементы $s_{i k}$ матрицы $S$ имеют вид:



\[

\begin{gathered}

s_{i k}(z)=z^{r_{i k}} \sum_{m=0}^{m_{i k}} \sigma_{i k m}(z) \ln ^{m} z, \\

\sigma_{i k m}(z)=\sum_{l=0}^{\infty} \sigma_{i k m l} z^{-l / n} .

\end{gathered}

\]



Эти ряды сходятся лишь в исключительных случаях и являются асимптотич. разложениями нек-рой фундаментальной матрицы в нек-рых секторах комплексной плоскости $z$ при $|z| \rightarrow \infty$. Асимптотика фундаментальной системы решений уравнения 2-го порядка



\[

w^{\prime \prime}-z^{r}\left(a_{0}+a_{1} z^{-1}+\ldots\right) w=0

\]



дается ВКБ-формулой



\[

w_{1,2} \sim z^{-r / 4} \exp \left( \pm \int_{z_{0}}^{z} t^{r / 2}\left(a_{0}+a_{1} t^{-1}+\ldots\right)^{1 / 2} d t\right)

\]



(см. Квазиклассическое приближение) при $|z| \rightarrow \infty$, $z$ лежит в ceкторе $\alpha<\arg z<\beta, \beta-\alpha<2 \pi /(r+2)$.





\end{document}